\section{尺度空间与细节空间}
\label{sec:wavelet-scaling-detail}

上一节的$2^J$个取值可以看成一个函数：在区间$[0,1)$的第$k$段上取常值$a_k$。这样，平均与差值就变成了对函数的运算，Haar变换就变成了在同一个函数空间中更换基函数。本节先建立由分段常数函数组成的尺度空间，再用正交补定义细节空间，最后核对取值向量与函数坐标之间相差的归一化因子。全节在$L^2([0,1])$中讨论，内积为$\langle f,g\rangle=\int_0^1\overline{f(x)}g(x)\,dx$，按第~\ref{chap:inner-product-spaces}章的约定对第一个变量共轭线性；本章出现的基函数都是实值的。

\subsection{尺度空间}
\label{sec:wavelet-scaling-spaces}

对整数$j\geq0$和$0\leq k<2^j$，记
\[
I_{j,k}=\bigl[k2^{-j},(k+1)2^{-j}\bigr),
\]
称为第$j$层的\term{二进区间}（dyadic interval）。第$j$层的$2^j$个二进区间互不相交，并起来是$[0,1)$。尺度$j$决定区间长度$2^{-j}$，位置$k$决定区间落在哪里。端点$x=1$只是一个零测集，在$L^2([0,1])$中不影响任何元素，所以下面不再区分$[0,1)$与$[0,1]$。

\begin{definition}[Haar尺度空间]
\label{def:wavelet-scaling-space}
对整数$j\geq0$，令$V_j$为$L^2([0,1])$中在每个$I_{j,k}$上几乎处处取常值的元素组成的子空间，称为第$j$层的\term{尺度空间}（scaling space）。令
\[
\varphi_{j,k}=2^{j/2}\mathbf{1}_{I_{j,k}}\qquad(0\leq k<2^j).
\]
\end{definition}

因子$2^{j/2}$的作用是归一化。不同的$\varphi_{j,k}$支撑在不相交的区间上，内积为零；而$\lVert\varphi_{j,k}\rVert_2^2=2^j\cdot2^{-j}=1$。$V_j$中的每个元素都是$\sum_kb_k\mathbf{1}_{I_{j,k}}$的形式，所以$\{\varphi_{j,k}\}_{0\leq k<2^j}$是$V_j$的标准正交基，$\dim V_j=2^j$。有限维子空间是闭的，第~\ref{chap:inner-product-spaces}章的正交投影可以直接使用。

\paragraph{嵌套}
每个第$j$层区间恰好是两个第$j+1$层区间的并：$I_{j,k}=I_{j+1,2k}\cup I_{j+1,2k+1}$。在$I_{j,k}$上取常值的函数，在这两个子区间上当然也取常值，所以
\[
V_0\subset V_1\subset V_2\subset\cdots .
\]
层级越高，允许的分段越细，空间越大。$V_0$只含常数函数。

\paragraph{投影就是取区间平均}
$L^2([0,1])$中的函数在有限测度区间上可积（由Cauchy--Schwarz不等式，$\int_I|f|\leq\lvert I\rvert^{1/2}\lVert f\rVert_2$），所以每个区间上的平均有意义。设$P_j$是到$V_j$的正交投影。用标准正交基写出，
\begin{equation}
P_jf=\sum_{k=0}^{2^j-1}\langle\varphi_{j,k},f\rangle\varphi_{j,k}
=\sum_{k=0}^{2^j-1}\Bigl(2^j\int_{I_{j,k}}f(x)\,dx\Bigr)\mathbf{1}_{I_{j,k}} .
\label{eq:wavelet-projection-average}
\end{equation}
括号中的数正是$f$在$I_{j,k}$上的平均。所以$P_jf$在每个第$j$层区间上取$f$在该区间的平均值，它是$V_j$中离$f$最近的函数。例如$f(x)=x$时，$P_0f\equiv1/2$，$P_1f$在$[0,1/2)$上为$1/4$、在$[1/2,1)$上为$3/4$。

\subsection{细节空间}
\label{sec:wavelet-detail-spaces}

从$V_j$到$V_{j+1}$，空间多出了$2^{j+1}-2^j=2^j$维。多出来的部分应当恰好记录“每个第$j$层区间的左右两半有何不同”，因为在第$j$层这一差别被平均掉了。用正交补把这部分单独取出。

\begin{definition}[Haar细节空间与Haar函数]
\label{def:wavelet-detail-space}
对整数$j\geq0$，令$W_j=V_{j+1}\cap V_j^{\perp}$，称为第$j$层的\term{细节空间}（detail space）。令
\[
\psi=\mathbf{1}_{[0,1/2)}-\mathbf{1}_{[1/2,1)},\qquad
\psi_{j,k}=2^{j/2}\bigl(\mathbf{1}_{I_{j+1,2k}}-\mathbf{1}_{I_{j+1,2k+1}}\bigr)\qquad(0\leq k<2^j),
\]
称$\psi_{j,k}$为\term{Haar函数}（Haar function）。
\end{definition}

$\psi_{j,k}$在$I_{j,k}$的左半取$2^{j/2}$、右半取$-2^{j/2}$，在其余位置为零。用变量代换可以写成$\psi_{j,k}(x)=2^{j/2}\psi(2^jx-k)$，所以全部Haar函数都由同一个$\psi$平移、缩放而来。

\begin{theorem}[细节空间的标准正交基]
\label{thm:wavelet-detail-basis}
对每个$j\geq0$，$\{\psi_{j,k}\}_{0\leq k<2^j}$是$W_j$的标准正交基，并且$V_{j+1}=V_j\oplus W_j$，两者正交。
\end{theorem}

\begin{proof}
$\psi_{j,k}$在每个第$j+1$层区间上取常值，属于$V_{j+1}$。它与$\varphi_{j,k'}$正交：$k'\neq k$时支撑不交；$k'=k$时
\[
\langle\varphi_{j,k},\psi_{j,k}\rangle=2^{j}\bigl(2^{-j-1}-2^{-j-1}\bigr)=0 .
\]
所以$\psi_{j,k}\in V_j^{\perp}$，从而$\psi_{j,k}\in W_j$。不同的$\psi_{j,k}$支撑不交，且$\lVert\psi_{j,k}\rVert_2^2=2^j\cdot2^{-j}=1$，它们是标准正交的。

有限维空间$V_{j+1}$中，闭子空间$V_j$的正交补$W_j$满足$V_{j+1}=V_j\oplus W_j$，维数为$2^{j+1}-2^j=2^j$。$W_j$中有$2^j$个标准正交向量，它们线性无关，因而构成$W_j$的基。
\end{proof}

也可以直接看出$V_{j+1}$中的每个函数都能用两组基表示：由定义，
\[
\mathbf{1}_{I_{j+1,2k}}=\frac{2^{-j/2}}{2}\bigl(\varphi_{j,k}+\psi_{j,k}\bigr),\qquad
\mathbf{1}_{I_{j+1,2k+1}}=\frac{2^{-j/2}}{2}\bigl(\varphi_{j,k}-\psi_{j,k}\bigr).
\]
这两个式子正是平均与半差的函数版本：一个细区间上的指示函数，等于所在粗区间的平均部分加上或减去一个细节。

\paragraph{逐层分解}
把定理~\ref{thm:wavelet-detail-basis}反复用于$V_J,V_{J-1},\ldots,V_1$，得到正交直和分解
\begin{equation}
V_J=V_0\oplus W_0\oplus W_1\oplus\cdots\oplus W_{J-1}.
\label{eq:wavelet-finite-decomposition}
\end{equation}
维数核对为$1+1+2+\cdots+2^{J-1}=2^J$。相应地，$\{\varphi_{0,0}\}\cup\{\psi_{j,k}:0\leq j<J,\ 0\leq k<2^j\}$是$V_J$的另一组标准正交基。

用投影写出，到$W_j$的正交投影是$P_{j+1}-P_j$。所以第$j$层细节就是“用第$j+1$层区间平均表示”与“用第$j$层区间平均表示”之差，正是较粗表示中缺少的局部差异。仍以$f(x)=x$为例，$P_1f-P_0f$在左半为$-1/4$、右半为$1/4$。另一方面$\langle\psi_{0,0},f\rangle=\int_0^{1/2}x\,dx-\int_{1/2}^1x\,dx=1/8-3/8=-1/4$，于是$\langle\psi_{0,0},f\rangle\psi_{0,0}$在两半上分别为$-1/4$和$1/4$，两种算法一致。

\subsection{值向量与函数坐标}
\label{sec:wavelet-values-coordinates}

现在把第~\ref{sec:wavelet-finite-haar}节的有限计算准确地放进$V_J$。设$f=\sum_{k=0}^{N-1}a_k\mathbf{1}_{I_{J,k}}$，$N=2^J$，即$f$在第$k$段上取值$a_k$。它的平方积分是
\[
\lVert f\rVert_2^2=\sum_{k=0}^{N-1}a_k^2\cdot2^{-J},
\]
与取值向量的欧氏平方和差一个因子$2^{-J}$。原因是每段的长度为$2^{-J}$。改用标准正交基表示，$f=\sum_k(2^{-J/2}a_k)\varphi_{J,k}$，所以$f$在基$\{\varphi_{J,k}\}$下的坐标是$2^{-J/2}\bm a$。这个坐标向量的欧氏长度才等于$\lVert f\rVert_2$。

\paragraph{一步变换就是一次换基}
在相邻两个细区间上，由定义有
\[
\varphi_{j,k}=\frac{\varphi_{j+1,2k}+\varphi_{j+1,2k+1}}{\sqrt2},\qquad
\psi_{j,k}=\frac{\varphi_{j+1,2k}-\varphi_{j+1,2k+1}}{\sqrt2}.
\]
若$f$在这两个细区间上的部分是$c_0\varphi_{j+1,2k}+c_1\varphi_{j+1,2k+1}$，代入上式验证可得
\[
c_0\varphi_{j+1,2k}+c_1\varphi_{j+1,2k+1}
=\frac{c_0+c_1}{\sqrt2}\varphi_{j,k}+\frac{c_0-c_1}{\sqrt2}\psi_{j,k}.
\]
右端的两个系数正是式~\eqref{eq:wavelet-h2}中$\bm H_2$作用在$(c_0,c_1)$上的结果。所以对坐标向量$2^{-J/2}\bm a$作有限Haar变换，得到的正是$f$在基$\{\varphi_{0,0}\}\cup\{\psi_{j,k}\}$下的坐标，排序与第~\ref{sec:wavelet-orthogonal-transform}小节的约定一致：先是$\langle\varphi_{0,0},f\rangle$，再是由粗到细的$\langle\psi_{j,k},f\rangle$。

\begin{example}[四段函数的Haar系数]
\label{ex:wavelet-four-values-function}
取$J=2$，$f$在四段上依次取$2,4,8,6$。坐标向量是$2^{-1}\bm a=(1,2,4,3)$。第一层得$(3/\sqrt2,-1/\sqrt2)$与$(7/\sqrt2,1/\sqrt2)$，第二层把$(3/\sqrt2,7/\sqrt2)$换成$(5,-2)$，所以
\[
\bigl(\langle\varphi_{0,0},f\rangle,\langle\psi_{0,0},f\rangle,\langle\psi_{1,0},f\rangle,\langle\psi_{1,1},f\rangle\bigr)=\Bigl(5,-2,-\frac{1}{\sqrt2},\frac{1}{\sqrt2}\Bigr).
\]
这些数也可以直接积分核对。$\langle\varphi_{0,0},f\rangle=\int_0^1f=(2+4+8+6)/4=5$是$f$的平均；$\langle\psi_{0,0},f\rangle=(2+4)/4-(8+6)/4=-2$；$\langle\psi_{1,0},f\rangle=\sqrt2(2/4-4/4)=-1/\sqrt2$；$\langle\psi_{1,1},f\rangle=\sqrt2(8/4-6/4)=1/\sqrt2$。

它们与例~\ref{ex:wavelet-four-values-orthogonal}的欧氏系数$(10,-4,-\sqrt2,\sqrt2)$逐项相差因子$2^{-J/2}=1/2$。能量核对为$\lVert f\rVert_2^2=(4+16+64+36)/4=30$与$25+4+1/2+1/2=30$。此时保持的是函数的平方积分，而不是取值向量的平方和。
\end{example}

这个归一化在实际计算中容易混淆：直接对取值向量作正交Haar变换，得到的数要再乘以$2^{-J/2}$才是$L^2$中的系数。若只关心系数的相对大小，因子无关紧要；若要用系数估计函数的$L^2$误差，就必须带上它。
